\chapter{Formal Methods as the Default Path}
\label{ch:vision}

\index{formal methods}
Formal methods are usually presented as an extra stage: first build a system, then
translate a simplified version into a model, and finally ask a specialist tool a
few questions. That workflow taxes correctness. The model drifts, the difficult
translation is postponed, and proofs concern yesterday's architecture.

Our aim is the reverse. The shortest path to a running service should pass through
an executable behavioural specification. Writing the specification should remove
work: it supplies message types, scenario tests, model-checking input,
documentation, and the skeleton of the service itself. Going without it should be
the inconvenient choice.

\section{Behaviour before machinery}

The primary object is a set of permitted, visible histories. An implementation may
use threads, actors, queues, databases, caches, or continuations. Those decisions
matter operationally, but they are not automatically part of the contract.
Promoting all internal state into the model creates a product state space whose
size grows exponentially with independently varying components. Worse, it makes
the specification brittle under harmless refactoring.

We therefore begin with \emph{behaviour}: who communicated with whom, what kind of
event occurred, what information it carried, and what causal history it extends.
Internal state enters only when it changes an observable promise.

\section{Three ingredients}

The method combines three traditions.

\begin{enumerate}[label=\arabic*.]
  \item Language theory gives compact generative descriptions of legal histories.
        Regular and context-free grammars provide the approachable starting point.
        Dana Richards and Henry Hamburger, the author's college professors, are
        cited in \citet{richards2017logic} as the source from which the author
        learned these methods, not as their originators. For the underlying habits of
        symbolization, validity, proof, quantification, and modal reasoning, we also
        draw on the accessible progression of \citet{gensler2017introduction}.
  \item Protocol analysis contributes explicit participants, adversarial choices,
        trace semantics, and tool-supported checking. The CSP treatment of
        \citet{ryan2001modelling}, including Steve Schneider, is cited as the source
        from which the author learned this treatment.
  \item Temporal logic turns expectations into queries over the generated
        transition system. CTL is central because it distinguishes claims about
        some possible future from claims about every possible future
        \citep{clarke1986automatic}.
\end{enumerate}

\section{A test for every abstraction}

Each abstraction in this book should answer a practical question:

\begin{quote}
Can a developer state a meaningful behaviour once, compose it with other
behaviours, check its consequences, and keep it stable while implementation details
change?
\end{quote}

If an abstraction does not improve that workflow, it does not belong in the core
library.
